\documentclass[11pt]{article}
\usepackage[T1]{fontenc}
\usepackage{textcomp}
\usepackage[margin=0.75in]{geometry}
\usepackage{amsmath,amssymb,amsthm}
\usepackage{array,booktabs}
\usepackage{hyperref}
\setlength{\parindent}{0pt}
\newtheorem{theorem}{Theorem}
\theoremstyle{definition}
\newtheorem{definition}{Definition}
\title{Flow Proof Artifact: prop-09-secant-tangent-converse}
\author{Generated by \texttt{flow doc proof}}
\date{Flow Proof Book}
\begin{document}
\maketitle
If from a point outside a circle two straight lines are drawn, one cutting the circle and one meeting it, and the rectangle on the whole and the exterior segment equals the square on the other, then the other is tangent.\\[0.5em]
\textbf{Source.} euclid — Elements, Book III, Proposition 9\\[0.5em]
\bigskip
\noindent{\large \textbf{Derived fact 1.} \textit{If from a point outside a circle two straight lines are drawn, one cutting the circle and one meeting it, and the rectangle on the whole and the exterior segment equals the square on the other, then the other is tangent} \hfill the Euclidean plane \cdot Euclid Book III \cdot Proposition 9: if the rectangle equals the square then the line is a tangent}\par
\smallskip\noindent\textit{Coordinate: the Euclidean plane \cdot Euclid Book III \cdot Proposition 9: if the rectangle equals the square then the line is a tangent}\par
\smallskip\noindent\textit{Source: euclid — Elements, Book III, Proposition 9}\par
\smallskip\noindent\textit{Built on: proposition 8: the rectangle on a secant and its exterior segment equals the square on the tangent, for Euclid Book III on the Euclidean plane}\par
\medskip
\noindent\textbf{Goal.} If from a point outside a circle two straight lines are drawn, one cutting the circle and one meeting it, and the rectangle on the whole and the exterior segment equals the square on the other, then the other is tangent\par
\noindent$\displaystyle \text{rect equals sq implies tangent}$\par
\medskip
\noindent
\begin{tabular}{@{} >{\raggedright\arraybackslash}p{0.06\textwidth} >{\raggedright\arraybackslash}p{0.47\textwidth} >{\raggedleft\arraybackslash}p{0.41\textwidth} @{}}
\textbf{\#} & \textbf{Proof} & \textbf{Mathematics} \\
\midrule
\textbf{1.} & We prove that proposition 9: if the rectangle equals the square then the line is a tangent for Euclid Book III the Euclidean plane. & \\[0.45em]
\textbf{2.} & Let P = exterior\_point. & \\[0.45em]
\textbf{3.} & Let secant = line\_PAB. & \\[0.45em]
\textbf{4.} & Let other = line\_PQ. & \\[0.45em]
\textbf{5.} & From step 2, step 3, and step 4, we can deduce that other is tangent. & $\displaystyle \text{other is tangent}$ \\[0.45em]
\textbf{6.} & We invoke the derived fact governing Euclid Book III the Euclidean plane: proposition 8: the rectangle on a secant and its exterior segment equals the square on the tangent, for Euclid Book III on the Euclidean plane. & \\[0.45em]
\textbf{7.} & From step 6, this implies rect equals sq implies tangent. Hence proven. & $\displaystyle \text{rect equals sq implies tangent}$ \\[0.45em]
\bottomrule
\end{tabular}
\label{thm:Geometry-Euclid Book III-proposition_9_if_the_rectangle_equals_the_square_then_the_line_is_a_tangent}
\par\medskip\hrule
\end{document}
